\section{Preliminaries And Evidence Boundary}
\label{sec:preliminaries}

For a prime stage head $p$, let
\begin{equation*}
M_p=\prod_{q\lt p}q
\end{equation*}

and retain the integers coprime to $M_p$. Their accepted residues repeat
modulo $M_p$; adjacent differences in one complete period form a cyclic
gap list. The verified companion article establishes accepted-value
completeness, strict increase, period shift, exact survivor count, and
the copy-or-merge transition rule.

This article takes those verified construction facts as inputs. It then
studies mathematical properties of 2-gap populations and signed interval
discrepancies. The division is deliberate: the companion article fixes
and machine-checks the sieve object; this article expands the
mathematical theory of that object without claiming that every new
theorem has a Scala encoding.

\subsection{Roles, Populations, And Scopes}
\label{sec:roles-populations-scopes}

The notation follows the shared research vocabulary:

\begin{itemize}
  \item $Q$ is a fixed \textbf{future head} used for square-safe
    certification;
  \item $r$ is one \textbf{incoming prime}, and $r_i$ is filter $i$ in a
    complete conditioned chain to $Q$;
  \item a \textbf{complete period} contains one full cyclic sieve pattern;
  \item a \textbf{local window} is an explicitly stated bounded integer
    interval;
  \item a \textbf{2-gap-start population} counts starts $x$, not all
    accepted values; and
  \item a \textbf{square-safe certificate} requires both $x$ and $x+2$ to
    lie strictly below $Q^2$ after every missing prime below $Q$ has been
    installed.
\end{itemize}

For a conditioned chain, $S_i$ denotes the local 2-gap starts immediately
before filter $r_i$, and $N_i=|S_i|$. The final survivor population is
$S_m$ with size $N_m$. Complete-period populations are named separately
and are never substituted for $N_i$ without an explicit localization
theorem.

\subsection{Evidence Status}
\label{sec:evidence-status}

The following status convention is used:

\begin{itemize}
  \item \textbf{Verified foundation:} supported by maintained Stainless
    source through the companion article.
  \item \textbf{Mathematically proved:} a complete mathematical proof is
    included here, but no corresponding \texttt{.holds} theorem
    currently exists.
  \item \textbf{Open:} the required estimate is stated but not proved.
\end{itemize}

These local labels map to the shared vocabulary as follows:
\textbf{Verified foundation} means a Stainless-verified theorem,
\textbf{Mathematically proved} means a mathematically proved theorem, and
\textbf{Open} covers an open hypothesis or an unproved required
estimate. The article's derivations and appendices are the authority for
its mathematical claims; linked property records are supplementary
navigation through the wider research program.

No theorem in this article claims infinitely many twin primes.

Each property section states its population, scope, and quantifier, then
proves the result. Verified construction inputs link to maintained Scala
contracts. A mathematical proof is not called Stainless-verified merely
because its finite instances can be computed.
